\chapter{Organizacja procesu zarządzania incydentami i plan reagowania}\label{ch:organizacja-procesu}

\section{Architektura procesu zarządzania incydentami}\label{sec:architektura-procesu}

\section{Główny plan reagowania na incydenty (IRP)}\label{sec:plan-reagowania}

\subsection{Faza przygotowania (Preparation)}\label{subsec:faza-przygotowania}

\subsection{Faza wykrywania i analizy (Detection and Analysis)}\label{subsec:faza-wykrywania}

\subsection{Faza ograniczania, eliminacji i odtwarzania (Containment, Eradication, Recovery)}\label{subsec:faza-ograniczania}

\subsection{Faza działań poincydentowych (Post-Incident Activity)}\label{subsec:faza-poincydentowa}

\section{Struktura organizacyjna, role i odpowiedzialności}\label{sec:role-odpowiedzialnosci}

\subsection{Zespół reagowania na incydenty (CSIRT / SOC) i interesariusze}\label{subsec:zespol-csirt}

\subsection{Macierz odpowiedzialności RACI}\label{subsec:macierz-raci}

\section{Klasyfikacja, priorytetyzacja i eskalacja incydentów}\label{sec:klasyfikacja-eskalacja}

\subsection{Kategorie i taksonomia incydentów bezpieczeństwa}\label{subsec:kategorie-incydentow}

\subsection{Macierz oceny wpływu, pilności i priorytetyzacji}\label{subsec:macierz-priorytetyzacji}

\subsection{Ścieżki eskalacji funkcjonalnej i hierarchicznej}\label{subsec:sciezki-eskalacji}

\section{Zasady komunikacji wewnętrznej i zewnętrznej}\label{sec:zasady-komunikacji}

\subsection{Komunikacja wewnętrzna i kanały awaryjne}\label{subsec:komunikacja-wewnetrzna}

\subsection{Raportowanie do organów nadzorczych i regulatorów (KNF, CSIRT, UODO)}\label{subsec:raportowanie-regulatorzy}

\subsection{Komunikacja kryzysowa z klientami i mediami}\label{subsec:komunikacja-klienci}
